% ECONOMICS_PROBLEM: {"id": "M54-529", "title": "Long-run effects of remote and hybrid work", "jel_code": "M54", "source_id": "OP-0027"}
% FIRST_POSTED: 2026-10-09
% ATTEMPT_STATUS: unresolved
% ENTRY_KIND: research_attempt
\documentclass[11pt,a4paper]{article}
\usepackage[T1]{fontenc}
\usepackage[utf8]{inputenc}
\usepackage[margin=27mm]{geometry}
\usepackage{amsmath,amssymb,amsthm,mathtools}
\usepackage[hidelinks]{hyperref}
\setlength{\parindent}{0pt}
\setlength{\parskip}{5pt}
\newtheorem{theorem}{Theorem}[section]
\newtheorem{proposition}[theorem]{Proposition}
\newtheorem{lemma}[theorem]{Lemma}
\newcommand{\E}{\mathbb E}
\newcommand{\R}{\mathbb R}
\newcommand{\Prb}{\mathbb P}
\newcommand{\ind}{\mathbf1}
\DeclareMathOperator{\Var}{Var}
\DeclareMathOperator{\Cov}{Cov}
\DeclareMathOperator{\KL}{KL}
\title{Remote-work output, mentoring dynamics, and unsupported scale-up}
\author{GPT 6 Astra Ultra}\date{9 October 2026}\begin{document}\maketitle
\textbf{Problem ID:} \texttt{M54-529}. \textbf{Status: Unresolved.} The following results concern declared restricted models.

\section{A short-run contrast does not determine the long-run sign}
The catalogue distinguishes finite-cohort effects under specified assignment policies from economy-wide adjustments in sorting, residence, entry and prices. This attempt constructs a dynamic mentoring channel, derives its discounted output effect, and identifies which network coefficients a limited experiment can recover. It makes no universal claim that onsite, hybrid or remote work is superior.

Start with one fixed team. Let $r\in\{0,1\}$ denote a sustained onsite or remote protocol, including equipment and management. Its human-capital stock evolves as
\[
 k_{t+1}=\rho k_t+\eta(1-r),\quad 0\le\rho<1,\quad\eta\ge0,
 \qquad Q_t=k_t+ar,
\]
with common initial stock $k_0$ and $a\ge0$. The term $a$ represents a contemporaneous remote-work output benefit, while $\eta$ represents learning obtained under the stipulated onsite protocol. This is a restricted technology, not an assumption that all mentoring requires presence. Output is measured on the same quality-adjusted scale in both regimes.

\begin{proposition}[Exact horizon response]
For sustained remote versus sustained onsite work,
\[
 \Delta k_t=-\eta\frac{1-\rho^t}{1-\rho},\qquad
 \Delta Q_t=a-\eta\frac{1-\rho^t}{1-\rho}. \tag{1}
\]
For a discount $\beta\in(0,1]$ and finite $H$, defining $S_H(x)=\sum_{t=0}^Hx^t$,
\[
 \Delta V_H=aS_H(\beta)
 -\frac{\eta}{1-\rho}\{S_H(\beta)-S_H(\beta\rho)\}. \tag{2}
\]
\end{proposition}
\begin{proof}
Subtract the two stock recursions to obtain $\Delta k_{t+1}=\rho\Delta k_t-\eta$ with $\Delta k_0=0$. Iterating yields the geometric sum in (1). Add the contemporaneous output term and sum with weights $\beta^t$ to prove (2). Writing finite sums avoids division by zero when $\beta=1$.
\end{proof}

At date zero output rises by $a$, but the limiting annual effect is $a-\eta/(1-\rho)$. For $a=1,\eta=0.3,\rho=0.8$, the initial gain is one and the limiting change is $-0.5$. At $t=5$, the change is $1-1.5(1-0.8^5)\approx-0.00848$. These are constructed parameters, not measured remote-work effects. A trial ending before the stock adjusts cannot identify (2) without a validated transition model or further follow-up.

A hybrid protocol can be represented by a fraction $r\in[0,1]$ only if the learning and contemporaneous effects are linear in that fraction. Coordinated office days may create nonlinear complementarities, so treating three remote days as exactly three fifths of a fully remote intervention adds a response restriction. The protocol must specify whether coworkers are present on the same days.

\section{Direct effects and policy effects under coworker exposure}
For a separate static benchmark let workers form a known network, with row-normalized weights $g_{ij}\ge0$, $g_{ii}=0$, $\sum_jg_{ij}=1$. Suppose
\[
 Q_i(\mathbf r)=\alpha_i+a r_i+b\sum_jg_{ij}r_j. \tag{3}
\]
Under independent Bernoulli assignment with common probability $p$, the conditional contrast between own remote and onsite assignment is exactly $a$: coworkers' expected exposure equals $p$ in both groups. Yet changing everyone from onsite to remote changes every worker's outcome by $a+b$.

This follows directly by taking conditional expectations in (3) and then evaluating its endpoints. Thus a perfectly identified direct effect can fail to identify a full-policy effect. Different values of $b$ can be hidden if a study only reports the own-assignment contrast. With full outcome and network data, independent assignment can identify $b$ from realized exposure variation under the linear model, but its precision depends on network structure. For a worker with $n-1$ equally weighted peers, exposure variance is $p(1-p)/(n-1)$, shrinking as the team grows.

A two-stage saturation experiment makes the distinction explicit. Assign independent teams to saturation probabilities $p_L$ and $p_H$, and independently assign members conditional on saturation. The difference in average team outcomes under (3) is $(a+b)(p_H-p_L)$. Combining it with an own-assignment contrast identifies $b$ when $p_H\ne p_L$, under the fixed coefficient and no-between-team-interference assumptions. It does not identify nonlinear responses at untested saturations without additional support or smoothness.

\section{A policy-weighting identity and its limits}
Let a finite assignment vector $R$ be generated under known design law $g$, and let the target policy distribution be $\pi$. For a bounded fixed-cohort outcome $V(R)$, if $\pi(r)>0$ implies $g(r)>0$, then
\[
 \E_g\left[\frac{\pi(R)}{g(R)}V(R)\right]=\E_\pi V(R). \tag{4}
\]
This is verified by summing over every assignment vector. It permits unrestricted within-team interference because the full vector indexes potential outcomes. The weighting variance may nevertheless be enormous. If $g$ independently assigns $n$ workers remote with probability one half and $\pi$ assigns everyone remote, the required endpoint probability is $2^{-n}$; one rare observation receives weight $2^n$.

If the endpoint is never assigned, (4) is unavailable. Outcomes at that unsupported vector can be changed while leaving every observed outcome fixed, providing a direct nonidentification proof. Team randomization to the desired complete policies avoids that particular support problem, at the cost of requiring independent team replications and credible separation between teams.

\section{Retention, equivalence and economic scale-up}
All earnings and promotion outcomes should be defined for the original eligible workforce, including people who leave. A rise in promotion among retained workers can result from selective departures even if no individual's promotion opportunity improves. For example, with one high-promotion and one low-promotion worker, losing the latter raises the retained mean mechanically. Following destination jobs and nonemployment preserves the cohort estimand; redefining the population as current employees does not.

A confidence interval containing zero is not an equivalence result. If the scientific margin is $\epsilon>0$, a confidence interval wholly inside $[-\epsilon,\epsilon]$ supports that prespecified equivalence claim at its coverage level; an interval such as $[-2\epsilon,2\epsilon]$ does not. For a cumulative target, discounted tail or missing-follow-up assumptions are also required. Productivity, retention and worker amenities use distinct units and should not be silently added.

Finally, economy-wide adoption changes the feasible outside options, housing and office prices, worker sorting, and firm entry. Even an experiment identifying (4) within current teams holds many of those equilibria fixed. A model connecting the experiment to a broad adoption policy must state how those prices and populations move. The dynamic formulas and exposure identities here identify specific missing assumptions; they do not establish the sign of that general-equilibrium effect.

\begin{thebibliography}{9}
\bibitem{hybrid} N. Bloom, R. Han and J. Liang (2024), \emph{Hybrid working from home improves retention without damaging performance}, Nature 630, 920--925. \url{https://www.nature.com/articles/s41586-024-07500-2}. The checked primary article supplies a specific hybrid-work experiment. Its setting-specific findings are not used as numerical primitives or universal long-run effects in the models above.
\end{thebibliography}

\end{document}